\subsection{
  Replies to comments received from Scott Oser - Feb 11, 2013 }

\begin{enumerate}

%........................................................................................................
\vspace{5 mm}
\item
{\em
  ``SIGNIFICANT COMMENT: Table 6: In the table of BANFF parameters and best-fit values, the numbers listed for the 
  CCQE axial mass and resonant axial mass do not match what is in v8 of the BANFF technical note (see TN-106 Table 12).  
  This immediately raised a concern in my mind that you may have used the wrong numbers as inputs.  
  I'm hoping instead that the table is just in error and that you used the correct values as given in TN-106v8..  
  Please clarify this confusion.''
}\\
\\
The two notes are quoting different numbers.
TN-106 is quoting the actual value of $M_A$, whereas TN-141 is
quoting a multiplicative factor describing how $M_A$ moves away from
nominal: $M_A$ $\rightarrow$ $x M_A$.
This muliplicative factor is calculated by: 
postfit value / NEUT nominal MC value,
i.e., $x_{MAQE} = 1.049 = 1.269 / 1.21$,
and $x_{MARES} = 1.011 = 1.22 / 1.21$ (note that there is a typo in v8
of TN-106).
This has been clarified in v4 of the note.

%........................................................................................................
\vspace{5 mm}
\item
{\em
  ``MINOR REQUEST FOR CLARIFICATION: Table 10: Could you briefly state the method used to derive the numbers in this table?  
  Is it the same method used to produce the "star plots" described in Section 4.4?''
}\\
\\
%The same method is used as in the calculation of Table 9, and Figs. 19-23: 
For each systematic parameter, 20k toy-MC experiments were generated, with randomized values of that systematic
parameter generated from a gaussian distribution\footnote{
Except for the spectral function systematic, which was, as always,
thrown uniformly between 0 and 1.}, and the 1 standard deviation of the total number of events was calculated.
%For pure normalisation systematics,
%the result is the same (except from very small statistical
%fluctuations) as the top row of the corresponding table in Appendix E.
%However, for systematics described by response functions, differences do occur.

%........................................................................................................
\vspace{5 mm}
\item
{\em
  ``COMMENT ON SECTION 5.6: I understand that the oscillation group has recommended that the VaLOR '2+N' 
  fit be T2K's primary result.  I am happy with that choice, and think it is the correct one.  However, 
  I am bewildered that you guys are still advocating for the '2+8' as our primary result even after you've 
  successfully implemented '2+N'.  Although you claim that the '2+8' method is simpler, the exact opposite is true: 
  it is never simple to explain why we would choose to present a statistically incomplete, and technically incorrect, 
  treatment (such as '2+8') when we have actually done a statistically correct one ('2+N'). Quite simply, the 2+N 
  method is statistically correct, and the 2+8 method is only acceptable to the extent that it gave the same result 
  as the correct method.  I am relieved that the oscillation group as a whole has selected the 2+N method as our primary result, 
  which at least means we don't really have to fight about this now, but had the '2+8' method been selected I would now be 
  loudly advocating for KIT or MUSIC as our primary result.  
  (Let me take a moment however to thank you all for having listened to our concerns about 
  this aspect of your analysis and having taken the time to do the 2+N method!)''
}\\
\\
Our opinion on the choice of the primary analysis is unchanged.
Lots of discussions took place during the last collaboration meeting, where various arguments were heard against, 
or in favour, of both `2+N' and `2+8' fits. We do not wish to re-iterate the arguments here.
In the end, it must be recognised that `statistical correctness' should be judged taking into account 
both the intrinsic uncertainty of the systematic error evaluation and
the accuracy limit of the analysis, due to the various approximations made.
We recognize that whether to fit for 8 or 48 systematics is a matter of taste, especially since the two
methods give nearly identical results. 
We produced both `2+8' and `2+N' results and, although our own preference is for the `2+8' one, 
we are satisfied with the selection of any of our results as the primary T2K one.

%........................................................................................................

\end{enumerate}
